% ===================== Part IV - Results =====================
% Draft. Figures from AfterLab4/p4_report_figs.m (figures/report/*.pdf).
% Adjust the \includegraphics paths to where the PDFs sit in the Overleaf project.
% Only observations here - the explanations belong in the Discussion.

\subsection{Results}

\paragraph{Measurement noise.}
\Cref{fig:p4_noise} compares the IMU elevation measurement with the helicopter lying on the ground and hovering at the linearization point.
On the ground the noise stays within about $\pm0.2^\circ$, while in hover it oscillates with an amplitude of about $15^\circ$.
The same holds for all five measurements: the standard deviation is between about 25 and 125 times larger in hover, which is roughly three to four orders of magnitude in variance.
The spectrum in the right panel shows that the hover noise is concentrated in a peak around $60\,\mathrm{Hz}$, with smaller peaks at higher frequencies, while the ground noise has no such peak.
The hover covariance also has a strong correlation between $\dot{p}$ and $e$ (correlation coefficient $0.98$), which is absent on the ground.
The hover covariance is used as $R_d$ in all following tests.

\begin{figure}[htbp]
    \centering
    \includegraphics[width=\textwidth]{figures/p4_noise.pdf}
    \caption{IMU elevation measurement on the ground and in hover. Left and middle: \SI{1}{s} of the measurement with the mean removed, note the different $y$-axes. Right: power spectral density over \SIrange{10}{40}{s} of the same runs.}
    \label{fig:p4_noise}
\end{figure}

\paragraph{$Q_d = 0$.}
\Cref{fig:p4_qd_zero} shows the run where the filter was started with $Q_d = 0$.
The diagonal of $\hat{P}$ decreases throughout the run, and at $t = 19\,\mathrm{s}$ $\hat{P}_{pp}$ is almost three orders of magnitude below the baseline value.
Up to about $t = 5\,\mathrm{s}$ the estimate follows the encoder.
After that the estimate $\hat{p}$ stays much closer to the reference than the encoder does, while the encoder pitch oscillates with growing amplitude and reaches about $100^\circ$ at $t \approx 17\,\mathrm{s}$.
The run was stopped at $t = 19\,\mathrm{s}$.

\begin{figure}[htbp]
    \centering
    \includegraphics[width=\textwidth]{figures/p4_qd_zero.pdf}
    \caption{$Q_d = 0$. Top: pitch from the IMU, the encoder and the Kalman estimate that is fed to the controller. Bottom: diagonal elements of $\hat{P}$, compared with the baseline run.}
    \label{fig:p4_qd_zero}
\end{figure}

\paragraph{$Q_d \to \infty$.}
\Cref{fig:p4_qd_inf} compares $Q_d = 10^6 I$ with the baseline over the two pitch steps.
The pitch that is flown (top panel) is almost the same in the two runs.
The estimate $\hat{p}$ (middle panel) is equal to the IMU measurement whenever new data arrives, so it carries the full measurement noise.
This noise is largest while the helicopter holds $p_c = -17^\circ$ between $t = 13$ and $18\,\mathrm{s}$.
The same noise appears in $V_d$ (bottom panel), where the standard deviation is about three to four times larger than in the baseline.
The logged $\hat{P}$ for the five measured states settles at exactly the diagonal of $R_d$.

\begin{figure}[htbp]
    \centering
    \includegraphics[width=\textwidth]{figures/p4_qd_inf.pdf}
    \caption{$Q_d \to \infty$ compared with the baseline over the two pitch steps. Top: pitch measured by the encoder. Middle: pitch estimate. Bottom: differential voltage $V_d$.}
    \label{fig:p4_qd_inf}
\end{figure}

\paragraph{Distribution of $Q_d$.}
\Cref{fig:p4_qd_split} compares the baseline with a high $q$ on the rates and a high $q$ on the angles.
With high $q$ on the rates, the pitch response and the error in $\hat{p}$ are close to the baseline, while the error in $\hat{\dot{p}}$ is a wider noise band than in the baseline.
With high $q$ on the angles, the pitch oscillates after each step and undershoots to about $-32^\circ$ after the step at $t = 13\,\mathrm{s}$.
In this run the error in $\hat{p}$ is a noise band of up to about $\pm10^\circ$, and $\hat{\dot{p}}$ is smaller than the encoder rate, which shows up as the slow oscillation in the bottom panel.
Of the three runs, the baseline has the smallest estimation error in both $\hat{p}$ and $\hat{\dot{p}}$.

\begin{figure}[htbp]
    \centering
    \includegraphics[width=\textwidth]{figures/p4_qd_split.pdf}
    \caption{Baseline, high $q$ on the rates and high $q$ on the angles. Top: pitch measured by the encoder. Middle and bottom: estimation error in $\hat{p}$ and $\hat{\dot{p}}$, relative to the encoder.}
    \label{fig:p4_qd_split}
\end{figure}

\paragraph{Disconnecting new data.}
In \cref{fig:p4_newdata} the new-data signal was disconnected twice during flight (shaded areas).
In both intervals the estimate $\hat{p}$ stays near $0^\circ$, while the encoder pitch moves away and reaches $36^\circ$ and $23^\circ$.
At the same time the plotted diagonal elements of $\hat{P}$ grow, from about $10^{-4}$ up to between $10^{-2}$ and $10^{-1}$.
When the signal is reconnected, $\hat{p}$ jumps to the encoder value, and $\hat{P}$ falls back to its previous level within about $0.05\,\mathrm{s}$.

\begin{figure}[htbp]
    \centering
    \includegraphics[width=\textwidth]{figures/p4_newdata.pdf}
    \caption{New-data signal disconnected twice during flight (shaded). Top: pitch reference, encoder and estimate. Bottom: selected diagonal elements of $\hat{P}$.}
    \label{fig:p4_newdata}
\end{figure}

\paragraph{Kalman filter and Luenberger observer.}
\Cref{fig:p4_vs_luenberger} compares the Kalman filter with baseline $Q_d$ and the Luenberger observer from Part III, flown with the same controller and reference profile.
The pitch responses lie on top of each other through the whole profile.
The estimation errors in $\hat{p}$ and $\hat{e}$ follow the same slow pattern for both estimators, but the Kalman error has more high-frequency noise on top.
The RMS error in $\hat{p}$ is the same for the two ($0.5^\circ$), while the RMS error in $\hat{\dot{p}}$ is lower for the Kalman filter ($3.0$ against $4.4\,\mathrm{deg/s}$).

A run with $Q_d$ computed by Van Loan's method, with the same intensity as the baseline, gave nearly the same RMS estimation error as the baseline for all states, and is therefore not plotted.

\begin{figure}[htbp]
    \centering
    \includegraphics[width=\textwidth]{figures/p4_vs_luenberger.pdf}
    \caption{Kalman filter (baseline $Q_d$) and Luenberger observer (Part III) with the same controller and reference. Top: pitch measured by the encoder. Middle and bottom: estimation error in $\hat{p}$ and $\hat{e}$, relative to the encoder.}
    \label{fig:p4_vs_luenberger}
\end{figure}
